\documentclass{article}
\usepackage{graphicx} % Required for inserting images
\usepackage{xcolor} % Must be loaded before soul
\usepackage{soul} % For highlighting text
\usepackage[utf8]{inputenc}
\usepackage{newunicodechar}
\usepackage{tikz}
\usetikzlibrary{positioning}
\usepackage[utf8]{inputenc}
\usepackage[T1]{fontenc}
\usepackage{fontspec}           % For Greek text support with XeLaTeX/LuaLaTeX
\usepackage{polyglossia}        % Multilingual support
\setmainlanguage{english}
\setotherlanguage{greek}

\usepackage{geometry}
\geometry{margin=1in}

\usepackage{csquotes}           % Block quotes
\usepackage{enumitem}           % Better lists
\usetikzlibrary{arrows.meta, positioning, shapes.geometric}

\usepackage{fancyvrb}           % Verbatim environments for ASCII diagrams
\usepackage{setspace}           % Line spacing



% --- Greek support (XeLaTeX) ---
\usepackage{fontspec}
\newfontfamily\greekfont{Gentium Plus} % good polytonic Greek; try Times New Roman if needed
\newcommand{\gk}[1]{{\greekfont #1}}

\newunicodechar{₁}{$_1$}

\usepackage[style=mla]{biblatex} % MLA Citation Style for Bibliography
\addbibresource{references.bib} %Add References page for in text citation
\usepackage[colorlinks=true, linkcolor=blue, urlcolor=blue, citecolor=blue]{hyperref} %enable hyperlinks to be colored and interactable
\usepackage{changepage} %create individual blocks of text that can be altered apart from general document.

\newcommand{\ra}{\ensuremath{\rightarrow}}

\newcommand{\inlinenote}[1]{\par\noindent
  \fcolorbox{orange}{orange!20}{\parbox{0.97\linewidth}{#1}}\par\noindent}

\begin{document}

\section{Being-in-Motion: The Ontological Horizon of Motion and Time}

\subsection{$A_0$ as First Actuality}

The actualization chain begins not with perception but with the ontological conditions perception presupposes. $A_0$ designates the chain's first actuality: the dyadic structure of the sensible object and the \textit{aisthētikon} in their conjoint actuality, the perceiver-perceptible pair from which the chain's opening motion ($M_0 \rightarrow A_1$) proceeds. This first actuality is itself situated within an encompassing horizon---motion (\textit{kinēsis}, \gk{κίνησις}) and time (\textit{chronos}, \gk{χρόνος}), the co-implicated structures whose prior actuality is the necessary condition for any chain-actualization to proceed. Before sensible objects can exercise their sensible powers, before sense-organs can hold first actuality as capacity, before the rational soul can actualize the counting through which temporal succession becomes articulated, motion and time must already be in act. The chain cannot begin with a static potentiality awaiting external awakening; it must begin with actualities that are themselves active---both the encompassing horizon of motion and time and, against this horizon, the dyadic actuality of $A_0$ itself. The section's notation is settled: $A_n$ names actualization stages (nodes), and $M_n \rightarrow A_{n+1}$ names the transitions between them (the motion-subscript indexes the source actuality; the terminus is given by the actuality at the arrow's head, after \textit{Physics} V.1, 224b7--9\footnote{``\ldots\ it is that to which rather than that from which the motion proceeds that gives its name to the change'' (\textit{Physics} V.1, 224b7--9).}). The chain $A_0 \rightarrow A_1 \rightarrow A_2 \rightarrow A_3 \rightarrow A_4$ runs from this first actuality through perception, \textit{phantasia} (\gk{φαντασία}), articulated desire, and the \textit{orektikon} (\gk{ὀρεκτικόν}) transition to overt action.

The architectural claim that motion and time are jointly the first actuality is grounded in Aristotle's doctrine of the priority of actuality over potentiality in three senses: in account (\textit{logos}, \gk{λόγος}), in time, and in substance (\textit{ousia}, \gk{οὐσία}). Aristotle defends this priority in \textit{Metaphysics} (IX.8): the potential is intelligible only through the actuality it is potential \textit{for} (1049b12--17); any individual potentiality presupposes a prior actuality capable of bringing it into play (1049b17--1050a3); the \textit{telos} (\gk{τέλος}) is more fully real than the process directed toward it (1050a4--16). The threefold priority licenses the chain's iterated structure: each completed actuality serves as the unmoved originator of the motion that produces the next, and Aristotle insists, in a passage Burke marks as pivotal, that ``one actuality always precedes another in time right back to the actuality of the eternal prime mover'' (Burke, \textit{Grammar}, p.~253)\footnote{Burke's gloss is a near-verbatim of Aristotle in \textit{Metaphysics} (IX.8, 1050b1--6): ``Obviously, therefore, the substance or form is actuality. From this argument it is obvious that actuality is prior in substance to potentiality; and as we have said, one actuality always precedes another in time right back to the actuality of the eternal prime mover.''}. The iteration is licensed at every node by the perpetual priority of actuality, and Burke isolates the formal dimension: the priority is not first a temporal claim but a formal one, in which ``man is `prior' to boy because man has already attained its complete form whereas boy has not'' (Burke, \textit{Grammar}, pp.~261--262). The formal priority licenses the iterated chain by which each completed actuality functions as both \textit{telos} for what preceded it and \textit{archē} (\gk{ἀρχή})\footnote{I retain the Greek \textit{archē} for what English commonly renders ``principle'' or ``origin'' because the Aristotelian sense is more determinate. Aristotle's canonical analysis (\textit{Metaphysics} V.1 [\gk{Δ}.1], 1012b34--1013a23) distinguishes six senses of \textit{archē}; the relevant sense for the chain is the fourth: ``That from which (not as an immanent part) a thing first arises, and from which the movement or the change naturally first proceeds, as a child comes from the father and the mother, and a fight from abusive language'' (1013a7--10). Aristotle's summary captures the common feature: ``It is common, then, to all to be the first point from which a thing either is or comes to be or is known'' (1013a17--19). In the chain architecture developed here, each completed actuality functions as \textit{archē} in precisely this fourth sense---the productive source from which the next motion proceeds---and the English ``principle''/``origin'' loses this productive-source determinacy.} for what follows. $A_0$ is the limit case: an actuality without prior actuality, since motion and time are co-eternal. Aristotle states the co-eternity explicitly: ``it is impossible that movement should either come into being or cease to be; for it must always have existed'' (\textit{Metaphysics} XII.6, 1071b6--8), and the same passage yokes motion and time: ``nor can time come into being or cease to be; for there could not be a before and an after if time did not exist. Movement also is continuous, then, in the sense in which time is; for time is either the same thing as movement or an attribute of movement'' (1071b8--11).

The section's architecture must be read against its twofold project. The first is to demonstrate \textit{phantasia} as the soul's unifying capacity---that through which past (memory), present (perception), and future (anticipation) are held together as a single cohesive temporal narrative, without which no extended deliberation, no rational action, no emotion, and no continuous self-understanding would be possible. The second is to fully articulate the ontological structure behind the actualization of \textit{orexis}---the perception-to-action sequence---in its two registers: synchronically, as the iterated chain that unfolds within a discrete event from first perceptual encounter to completed action; and diachronically, as the cumulative trajectory within being-in-the-world, where each completed action reshapes the dispositional ground (\textit{hexeis} chief among them) from which subsequent actualizations proceed. Because this two-register articulation is the ontological structure of human being, it can serve equally as a methodology for analyzing any specific event from initial perception of stimuli through to completed action, or at minimum as a means of individuating the constitutive processes at work. The complete actualization chain the section develops is at once the perception-to-movement chain of \textit{De Motu Animalium} --- \textit{aisthēsis} or \textit{noēsis} or \textit{phantasia} $\rightarrow$ \textit{orexis} $\rightarrow$ \textit{pathē}\footnote{The Forster (Smith/Barnes ROT) translation of \textit{De Motu Animalium} 702a17--19 renders \textit{pathē} (\gk{πάθη}) as ``affections''; per the dissertation's Greek-throughout convention I retain \textit{pathē}. The chain-element corresponds to the basic-affective-valence and higher-order \textit{pathē} registers developed at \S 1.4.} $\rightarrow$ organic parts $\rightarrow$ movement (702a17--19) --- and the actualization chain of \textit{orexis} in which the source of all animal locomotion is articulated. These are not two parallel chains. They are one chain. Every cognitive act --- including pure intellection --- is itself a \textit{kinēsis} of the soul, since ``being pained or pleased, or thinking'' are themselves modes of motion of the ensouled being (\textit{De Anima} I.4, 408b5--7); and \textit{orexis} is the moved mover that produces all such motion: ``that which at once moves and is moved'' (\textit{De Anima} III.10, 433b17--18). The chain $A_0 \rightarrow A_4$ thus names one structure: $A_0$ is the dyadic structure of the sensible object and the \textit{aisthētikon}, situated within motion and time as the encompassing horizon within which all chain-actualizations unfold; the successive stages develop the single chain through which \textit{orexis} moves the animal via perception, \textit{phantasia}, evaluative articulation, and bodily preparation.\footnote{Struever's reading of Heidegger's lectures emphasizes precisely this ontological horizon as the \textit{Da-Charakter} of Dasein: ``\textit{kinesis} gives the authentic \textit{Da-Charakter} to Dasein, the vital place-time of Being.... the ontological significance of the passions lies in our capacity for change, the \textit{Veränderlichkeit} of Being'' (Struever, ``\textit{Alltäglichkeit}, Timefulness, in Heidegger's Early Lectures,'' in \textit{Heidegger and Rhetoric}, ed. Gross \& Kemmann, SUNY, 2005, pp.~109--110, citing GA 18, p.~287). Struever's broader claim that ``timefulness colonizes \textit{Being and Time}'s temporal program from within'' (p.~117) further supports motion and time as the ontological horizon: the horizon is not external to Dasein's existential structure but is what Dasein primarily \textit{is}, in its capacity for being-altered, being-moved, being-affected.}

The Heideggerian register operates from the section's opening, not adjacent to it. Heidegger glosses Aristotelian motion in the \textit{Basic Concepts of Aristotelian Philosophy} as the structural-existential determination of a being's \textit{there}: if a being is determined in its \textit{there} as Aristotle determines \textit{kinēsis}, then it is in movement; \textit{kinēsis} accordingly names what Heidegger develops as being-moved (\textit{Bewegtheit})\footnote{I follow Heidegger's 1924 reading in the \textit{Basic Concepts of Aristotelian Philosophy} (\textit{GA} 18, §§25--26), which reinterprets Aristotle's \textit{kinēsis} as \gk{ἐντελέχεια τοῦ δυνάμει ὄντος} (the actuality of what exists potentially, as such; \textit{Physics} III.1, 201a10): not ``motion'' in the modern physical sense (e.g., ``movement as uniform speed: c = s/t,'' which Heidegger contrasts at BCAP p.~199), but the presence of a being in its ability-to-be, precisely insofar as it is an ability-to-be. This is what Heidegger throughout calls being-moved (\textit{Bewegtheit}). The identification \textit{kinēsis} $\cong$ \textit{Bewegtheit} (in preference to the more conventional German rendering \textit{Bewegung}, ``movement'') is Heidegger's terminological move; I take it over as a framing-level commitment. On this reading the 1924 lectures anticipate the existential analytic of \textit{Being and Time}: \textit{kinēsis} is shown not as a regional natural-philosophical phenomenon but as ``the basic character of beings as moved'' (BCAP p.~199)---the mode of being of worldly beings whose way of being is to be underway-toward what they are not yet.}, the way of being of beings whose mode of existence is to be underway toward what they are not yet. The priority of actuality, on this reading, is not a free-standing metaphysical thesis about substance and capacity but a determination of how a comporting being holds itself open toward its ownmost possibilities --- including the possibilities that, at $A_0$, the entire chain of perception-to-action presupposes. The plan of the section follows from this dual register. Motion as \textit{energeia ateles} (\gk{ἐνέργεια ἀτελής} `incomplete actuality, actuality whose end lies beyond it') opens the analysis; the chain structure inherited from priority follows; the relational dyadic structure of \textit{kinēsis} and time as the number of motion develop the analysis through to the soul's constitutive role, at which point the threshold of $A_1$ is reached. 

\subsection{\textit{Kinēsis} as \textit{Energeia Ateles}: Being-in-Motion-and-Time}

Aristotle's compact definition of motion establishes the term whose ontological texture this subsection must specify:
\begin{adjustwidth}{0.5in}{0in}
the fulfilment of what is potentially, as such, is motion --- e.g. the fulfilment of what is alterable, as alterable, is alteration; of what is increasable and its opposite, decreasable (there is no common name for both), increase and decrease; of what can come to be and pass away, coming to be and passing away; of what can be carried along, locomotion. (\textit{Physics} III.1, 201a11--14)
\end{adjustwidth}
Motion is the actuality (\textit{energeia}, \gk{ἐνέργεια}) of the potential (\textit{dynamis}, \gk{δύναμις}) qua potential; it is a genuine actuality, and yet it is incomplete because its \textit{telos} lies outside itself. Aristotle's pairing in \textit{Metaphysics} sharpens the point. Some activities have a same-time structure: ``at the same time we are seeing and have seen, are understanding and have understood, are thinking and have thought'' (\textit{Metaphysics} IX.6, 1048b23--25). Others do not: ``it is not true that at the same time we are learning and have learnt, or are being cured and have been cured'' (1048b25--27). The structural difference is decisive: ``every movement is incomplete---making thin, learning, walking, building; these are movements, and incomplete movements'' (1048b28--30). Seeing and thinking are \textit{energeiai} in the strict sense, since they contain their \textit{telos} within themselves; the moment one has seen, one is seeing, and the activity does not cease on reaching its terminus. Building, learning, and curing are \textit{kinēseis} because their terminus lies beyond the activity itself; the moment the house exists, building has ceased.

The peculiar ontological texture of motion follows from this asymmetry. Motion is the kind of actuality that is at once genuinely present --- the builder is at work, the house is being-built --- and incomplete in respect of what it is becoming. Hawhee observes of this register that \textit{energeia} ``names something like physical presence: \textit{energeia} means the presence of the thing'' (Hawhee, ``Rhetorical Vision,'' p.~154). If \textit{energeia} names the presence of the thing, then \textit{kinēsis} as \textit{energeia ateles} names a peculiar mode of presence: the presence of something that is not yet fully what it is becoming. Aristotle's bronze statue example clarifies this. The motion of bronze becoming a statue is not the fulfilment of bronze \textit{as bronze} but the fulfilment of bronze \textit{as potentially a statue} (\textit{Physics} III.1, 201a28--33). The as-structure governs everything: what is actualized in motion is actualized \textit{as} potential, not as achieved completion. Aristotle distributes motion across four categories --- substance, quality, quantity, and place --- and finds in each a pair of contraries between which the moving thing travels (\textit{Physics} V.2, 226a23--30); the world is always already in motion, and there is no moment at which potentiality sits inert, waiting to be actualized from without. $A_0$ is always already actual, and it is actual as ongoing.

The distinction between the two-modes-of-being-affected Aristotle draws in \textit{De Anima} reinforces the point and prepares the relational analysis to come. ``The expression `to be acted upon' has more than one meaning; it may mean either the extinction of one of two contraries by the other, or the maintenance of what is potential by the agency of what is actual and already like what is acted upon, as actual to potential'' (\textit{De Anima} II.5, 417b2--5). In one sense being-affected is a destruction; in the other it is a preservation in which a potential is brought into its actuality. Hearing a sound is a being-affected of the second sort: my hearing is brought to its full being, not destroyed. The same structural point governs the motion-time horizon within which $A_0$ stands: it is not the kind of actuality that is exhausted in producing the next stage of the chain but the kind that preserves itself as the ground on which the next stage stands.

Heidegger's reading of Aristotelian \textit{kinēsis} in the 1924 Marburg lectures (\textit{BCAP}) transposes this analysis from the level of natural-philosophical category to the level of fundamental ontology. Motion, on Heidegger's reading, is a character of being (\textit{Seinscharakter}): if a being is determined in its \textit{there} as Aristotle determines \textit{kinēsis}, then that being is in movement, and \textit{kinēsis} accordingly names what Heidegger develops as being-moved (\textit{Bewegtheit}): ``Here, it is really not a question of \textit{defining} movement in some sense, but of \textit{making} beings as moved \textit{visible} in their being-there and holding fast to them'' (\textit{Basic Concepts of Aristotelian Philosophy}, p.~199).\footnote{BCAP §26 (pp.~195--212) develops Heidegger's reading of Aristotelian \textit{kinēsis} as movedness (\textit{Bewegtheit}). The chapter-title---``Movement as \gk{ἐντελέχεια τοῦ δυνάμει ὄντος}'' (the actuality of what exists potentially, as such; cf.~\textit{Physics} III.1, 201a10)---marks the reading-strategy: Heidegger reads \textit{Metaphysics} $\Theta$ on \textit{dynamis} and \textit{energeia} alongside \textit{Physics} $\Gamma$.1, recovering \textit{energeia ateles} at the level of fundamental ontology. Heidegger's analysis of \textit{dynamis} (BCAP §26d.γ) specifies that the possibility a being has is not arbitrary: ``Only what is cold has the possibility of the warm, not what is hard or red. \ldots\ The possibility is not just any arbitrary one, but rather one that has a definite direction'' (BCAP p.~196). This directedness of possibility is the structural feature motion presupposes, and it sets up the \textit{energeia}/\textit{dynamis} pair as a single concept-pair through which motion's being-character becomes intelligible.} The structural payoff is precise. Dasein --- the being whose way of being is to comport itself toward its own being --- is not a static substance which subsequently enters into changes; it is a being whose mode of existence is being-underway-toward, the ahead-of-itself structure that \textit{Being and Time} develops as care (\textit{Sorge}; \S 41). ``Dasein is its possibility,'' Heidegger writes (BT \S 9, H.42; Eng.~67), and the formula recovers Aristotelian motion at the level of fundamental ontology --- not as an external addition to a prior substance but as the way in which the comporting being is what it is. The intermediate-register character of \textit{energeia ateles} --- genuinely actual yet structurally incomplete --- is therefore not a marginal metaphysical curiosity but the primitive ontological pattern from which the entire chain inherits its iterated structure: each completed actuality at $A_1$, $A_2$, $A_3$ is both the terminus of its antecedent motion and the unmoved originator of its successor, an \textit{entelecheia} (\gk{ἐντελέχεια} `having-one's-end-within-oneself') in relation to what precedes it and an \textit{archē} in relation to what follows. The chain is iterated because each of its nodes has the same dual structure $A_0$ establishes as primitive. The temporality of taking care that \textit{Being and Time} develops --- the unity of awaiting, making-present, and retaining as the temporal articulation of Dasein's circumspective concern (BT \S 68b, H.342) --- recovers, at the existential level, what the chain establishes at the level of categorial structure: the comporting being is constituted by the ongoing motion of its being-underway, and each completed phase of its concern is at once the terminus of the motion that produced it and the principle of the motion that follows.

\subsection{Priority of Actuality and the Chain Structure}

The chain structure of $A_0 \rightarrow A_4$ inherits its logic from the priority of actuality over potentiality. Aristotle states the priority in its most concentrated form:
\begin{adjustwidth}{0.5in}{0in}
everything that comes to be moves towards a principle, i.e. an end. For that for the sake of which a thing is, is its principle, and the becoming is for the sake of the end; and the actuality is the end, and it is for the sake of this that the potentiality is acquired. For animals do not see in order that they may have sight, but they have sight that they may see. (\textit{Metaphysics} IX.8, 1050a7--11)
\end{adjustwidth}
The end is not merely the terminus of a process but its \textit{archē}: the completed form furnishes the reason why the process occurs at all and the measure by which its stages are ordered. Sight is not a capacity that becomes actualized as seeing; seeing is the actuality for the sake of which sight exists as a capacity. The three senses of priority each carry architectural consequences for the chain. In account, the potential is intelligible only through the actuality it is potential for; the chain's apparatus is therefore intelligible only through the completed $A_4$ that gives it its shape. In time, any individual potentiality presupposes a prior actuality capable of bringing it into play; within motion and time as the encompassing horizon, $A_0$'s dyadic actuality of sensible object and \textit{aisthētikon} is what brings into play the joint actualization at $A_1$'s dynamis-side. In substance, the \textit{telos} is more fully real than the process directed toward it; the completed action at $A_4$ is more fully real than the chain's earlier phases, even though those phases must precede it in time. Burke isolates the architectural payoff: the entelechy ``allowed [Aristotle] to introduce another kind of priority, namely the `principle' involved in a given form'' (Burke, \textit{Grammar}, pp.~261--262). The priority licenses an iterated structure in which each completed actuality functions as both terminus and origin.

The principle by which motions are named in the chain is articulated cleanly in the \textit{Physics}: ``it is that to which rather than that from which the motion proceeds that gives its name to the change'' (\textit{Physics} V.1, 224b7--8). Each $M_n \rightarrow A_{n+1}$ in the chain takes its name from its terminus, not from its starting point. The motion of perception is named for the perceptual \textit{energeia} it produces; the motion of \textit{phantasia}-formation, for the \textit{phantasma} (\gk{φάντασμα} `image, appearance') it leaves behind; the \textit{orektikon} transition, for the desire's discharge into action. Aristotle's geometer illustration sharpens the relation between potential and actual:
\begin{adjustwidth}{0.5in}{0in}
the potentially existing constructions are discovered by being brought to actuality (the reason being that thinking is the actuality of thought); so that the potentiality is discovered from actuality (and therefore it is by an act of construction that people acquire the knowledge). (\textit{Metaphysics} IX.9, 1051a29--33)
\end{adjustwidth}
All the possible geometric configurations exist potentially in the figure, but they are not discovered by introspection; the geometer who proves the theorem actualizes the potentialities that were always-already there yet latently, and the actualization makes them manifest. The principle generalizes: each stage's potentialities are not abstract capacities awaiting filling-in but structural possibilities that become intelligible only when the prior actuality has licensed them.

The transitional character of motion as encompassing horizon is articulated in Aristotle's general formulation of motion as passage between states: ``if there is movement and something moved, that which is moved must first be in that out of which it is to be moved, and then not be in it, and move into the other and come to be in it'' (\textit{Metaphysics} XI.6, 1063a19--21). Motion is the medium of passage; the chain's iterated structure is a sequence of such passages, each instantiating the same transitional ontology. The relation between the chain's stages is therefore not a causal-mechanical sequence but a hermeneutical one. The completed actuality at each stage --- the for-the-sake-of-which --- is always already operative in the understanding that guides the actualization. Aristotle articulates this in his own register as the operation of the \textit{telos}; Heidegger generalizes the structure at the level of Dasein's interpretive comportment as the hermeneutical circle. In \textit{Being and Time}, Heidegger insists: ``In every case this interpretation is grounded in \textit{something we have in advance}---in a \textit{fore-having}'' (BT §32, H.150; Eng.~191). The fore-structure has a triadic articulation---fore-having (\textit{Vorhabe}), fore-sight (\textit{Vorsicht}), and fore-conception (\textit{Vorgriff})---and Heidegger summarizes: ``Whenever something is interpreted as something, the interpretation will be founded essentially upon fore-having, fore-sight, and fore-conception'' (BT §32, H.150; Eng.~191). This triadic fore-structure articulates the as-structure of interpretation, in which what is interpreted is already disclosed as something prior to any propositional articulation. The circle is the structural condition for any actualization that proceeds toward a \textit{telos}: the \textit{telos} must already be operative in the comportment that aims at it, and the comportment must already be guided by the \textit{telos} it does not yet possess. ``What is decisive is not to get out of the circle but to come into it in the right way,'' Heidegger insists; ``this circle of understanding is not an orbit in which any random kind of knowledge may move; it is the expression of the existential fore-structure of Dasein itself'' (BT §32, H.153; Eng.~195). The chain's iterated structure --- \textit{energeia ateles} generating \textit{entelecheia} generating new conditions for \textit{energeia ateles} --- is the proto-form of the hermeneutical circle Heidegger generalizes for Dasein's interpretive comportment as a whole.\footnote{Heidegger's hermeneutical circle (BT §§32--33, H.150--160) and the iterative completion-structure of Aristotelian motion share the same projection-toward-completion schema. \textit{Energeia ateles} is the actuality of the potential \textit{qua} potential---a genuine actuality that is structurally incomplete because its \textit{telos} lies beyond itself yet is operative within the actualization as its directing principle (\textit{Physics} III.1, 201a10--11). On the Heideggerian side, understanding projects upon possibilities that are themselves disclosed by the fore-structure of understanding---the projection-toward-completion is already disclosed in the projection itself (BT §32, H.150). In both registers, the end is already operative in the comportment that proceeds toward it, and the comportment is intelligible only through the end it does not yet possess. Heidegger does not draw the parallel to Aristotle explicitly in BT §§32--33, but the projection-schemas operate at the existential and categorial-structural levels respectively.} Aristotle's apparatus articulates the structural-categorial features of the chain; Heidegger's apparatus articulates how those features are determinations of a comporting being's way of being. The section deploys the two registers together.\footnote{Hans-Georg Gadamer, in his first-person testimony from the SS 1923-24 Aristotle seminars, reports that Heidegger's reading of Aristotelian \textit{kinēsis} as \textit{Bewegtheit} (``movement'') and the corresponding ontology of being-underway-toward emerged through the seminar reading of \textit{Physics} III, \textit{Metaphysics} IX, and the \textit{Rhetoric}. Gadamer's formula --- ``It is the same \textit{dynamis}!'' --- registers the recognition that the rhetorical, hermeneutic, and natural-philosophical strands of Aristotle's text all converge on a single ontological structure (Gadamer 60). Michalski reads GA 18 along the same philological-philosophical seam: Heidegger ``rejects the title of philosophy in favor of philology'' (Michalski 72) precisely because the work of reading Aristotle's \textit{kinēsis}-doctrine \textit{is} the work of disclosing the existential-temporal structure the doctrine articulates. The two-register deployment the present section undertakes is, on this reading, the same gesture Heidegger first made in 1924.}

\subsection{The Relational Structure of \textit{Kinēsis}}

A feature of $A_0$ that the chain inherits at every subsequent stage is the relational character of \textit{kinēsis}. Motion is not a property inhering in a single substance but a relation between mover and moved, between the active and passive powers whose joint actualization constitutes the process of change. Aristotle states the relation directly:
\begin{adjustwidth}{0.5in}{0in}
motion is in the movable. It is the fulfilment of this potentiality by the action of that which has the power of causing motion; and the actuality of that which has the power of causing motion is not other than the actuality of the movable; for it must be the fulfilment of \textit{both}. A thing is capable of causing motion because it \textit{can} do this, it is a mover because it actually \textit{does} it. But it is on the movable that it is capable of acting. Hence there is a single actuality of both alike. (\textit{Physics} III.2, 202a13--20)
\end{adjustwidth}
Motion is irreducibly relational, a joint actualization of two capacities --- the capacity to act and the capacity to be acted upon --- and the \textit{energeia} shared between them is not the consequence of the motion but the motion itself. Aristotle specifies the dyadic structure in categorial terms: ``it is necessary that there should be an actuality of the agent and of the patient. The one is agency and the other patiency; and the outcome and end of the one is an action, that of the other a passion'' (\textit{Physics} III.3, 202a21--24). The agent's actuality is \textit{poiēsis} (\gk{ποίησις}); the patient's actuality is \textit{pathēsis} (\gk{πάθησις}). Aristotle clarifies the point with the Thebes-Athens analogy: ``Thebes and Athens are different cities, yet the road from Thebes to Athens and the road from Athens to Thebes are the same'' (\textit{Physics} III.3, 202b13--14). One road, two descriptions. One motion, two formal aspects. The distinction is in account, not in substrate.

The dyadic architecture supplied by $A_0$ propagates through every subsequent stage of the chain. In the $M_0 \rightarrow A_1$ motion from $A_0$'s dyadic actuality, the sensible object (agent) and the sense-organ (patient) jointly actualize in a single perceptual act issuing in $A_1$ as completed perception: perception is itself a kind of alteration, in which the sense-organ is moved by the sensible object, and the perceptual act is the joint actualization of the organ's capacity to receive form and the object's capacity to act upon it.\footnote{Aristotle develops the perceptual case of the agent/patient structure (\textit{Physics} III.3, quoted in the block-quote above) in \textit{De Anima} (II.5, 416b33--417a20): perception is the joint actualization of the sensible object's active potency (agency) and the sense-organ's receptive potency (patiency). The single-event-under-two-\textit{logoi} structure is given in \textit{De Anima} (III.2, 425b26--27): ``the actualization of the sensible object and that of the percipient sense is one and the same actualization, though their being is not the same.'' The two-\textit{logoi} structure is the basis for the cognition/affect articulation analyzed in \S 1.0.} The full analysis of this propagation belongs to the section on \textit{aisthēsis} (\gk{αἴσθησις}), but the structural inheritance is fixed at $A_0$. The same schema iterates at the cognitive-motor stages of the chain: what the realizable good moves at $M_3 \rightarrow A_4$ is \textit{orexis}, the moved mover that is at once agent (of the animal's locomotion) and patient (of the apparent good); what \textit{orexis} in turn moves is the animal via its bodily instrument. The chain is dyadic at every node because $A_0$ establishes dyadic relationality as the primitive structural feature of \textit{kinēsis}.

Heidegger's existential radicalization of the relational structure is the analysis of Dasein as constituted by its being-in-the-world (\textit{In-der-Welt-sein}). Dasein does not first exist as substance and then enter into relations; it is constituted by relational being-in-the-world, in which entities are encountered as ready-to-hand (\textit{Zuhandenheit}), as members of equipmental contexts (\textit{Zeug}) articulated by chains of involvement (\textit{Bewandtnis}). Heidegger states the structural claim in \textit{Being and Time}: ``The character of Being which belongs to the ready-to-hand is just such an involvement'' (BT \S 18, H.84; Eng.~115).\footnote{Heidegger develops involvement (\textit{Bewandtnis}) as the assignment-structure characteristic of ready-to-hand entities: ``If something has an involvement, this implies letting it be involved in something. The relationship of the `with \ldots in \ldots ' shall be indicated by the term `assignment' or `reference'\,'' (BT §18, H.84; Eng.~115). The totality of involvements (\textit{Bewandtnisganzheit}) is structurally prior to any individual item of equipment: ``the totality of involvements which is constitutive for the ready-to-hand in its readiness-to-hand, is `earlier' than any single item of equipment'' (BT §18, H.84; Eng.~116). The structural payoff for the Aristotelian chain: just as the dyadic mover-moved relationality \textit{is} the motion (rather than something added to two pre-given substances), the totality of involvements \textit{is} the worldhood of the world (rather than an aggregate of independent substances within it). The chain's iterated dyadic structure recovers, at the existential level, the same structural primacy of relation over relata.} The Aristotelian relational \textit{kinēsis} is, on this reading, the proto-form of Heideggerian worldly involvement:\footnote{The identification of the dyadic potentiality-actuality (\textit{dynamis}-\textit{energeia}) structure of Aristotelian \textit{kinēsis} (mover-moved; capacity-exercise) with the proto-form of Heideggerian being-in-the-world (\textit{In-der-Welt-sein}) is grounded in the structural homology between joint-actualization and unitary-phenomenon. Heidegger insists that ``the compound expression `Being-in-the-world' \ldots\ stands for a unitary phenomenon. This primary datum must be seen as a whole'' and ``cannot be broken up into contents which may be pieced together'' (BT §12, H.53; Eng.~78). Just as \textit{kinēsis} is the joint actualization of two co-constituted capacities rather than a property attaching first to a separable substance, so being-in-the-world is the unitary phenomenon out of which Dasein and world are intelligible together rather than a relation between a subject and a separate object-world. Heidegger's BCAP §26 anticipates this reading (\textit{kinēsis} as the fundamental mode of being of worldly beings, before the 1927 \textit{Zuhandenheit} analysis); the being-in-the-world structure is then developed as the worldhood-environment-circumspection complex at BT §§14--18, H.63--88.} motion is always a being-toward, an articulation between an active power and a passive power that are not first independent substances and then jointly active but are co-constituted in the actuality of the motion they share. The chain's iterated dyadic structure is, at the existential level, the structural-categorial trace of the way a being whose mode of existence is being-in-the-world is constituted by its relations to what it encounters.

\subsection{Time (\textit{Chronos}) as the Number of Motion}

Aristotle's celebrated definition of time establishes the second moment of the encompassing horizon within which $A_0$ stands:
\begin{adjustwidth}{0.5in}{0in}
time is the number of motion in respect of the `before' and `after'. (\textit{Physics} IV.11, 219b1--2)
\end{adjustwidth}
Time is not an independent substance --- not a container within which events transpire, not an autonomous metaphysical entity --- but an aspect of motion: motion insofar as it admits of enumeration. Where there is no motion there is no time, and where motion can be numbered there is time as that numbering. The continuity and directional structure of time are inherited from motion, and motion inherits these from spatial magnitude: ``the movement goes with the magnitude, and time, as we maintain, goes with motion'' (\textit{Physics} IV.11, 219a10--13). Time is accordingly ``either the same thing as movement or an attribute of movement'' (\textit{Metaphysics} XII.6, 1071b10--11).

The definition turns on a distinction that is easy to miss but architecturally load-bearing. ``Time, then, is a kind of number. Number, we must note, is used in two ways --- both of what is counted or countable and also of that with which we count. Time, then, is what is counted, not that with which we count: these are different kinds of things'' (\textit{Physics} IV.11, 219b5--10). Time belongs to the first category: it is the counted number of motion, not the counting activity. The distinction will prove decisive when the soul's role enters the analysis, for it allows time to be located in the structure of motion itself --- as what \textit{can} be counted --- while still requiring a counter for the full actualization of time as numbered. The grounding of the definition is at once structural and experiential. ``We perceive movement and time together,'' Aristotle observes, ``even when it is dark and we are not being affected through the body, if any movement takes place in the mind we at once suppose that some time has indeed elapsed'' (\textit{Physics} IV.11, 219a3--10). Neither motion nor time is epistemically prior to the other; they are apprehended together, and the reciprocity is ontological as well as epistemic.

The `now' (\gk{τὸ νῦν} `the now') is the principle by which time's continuum is at once articulated and divided. ``The `now' determines time, insofar as time involves the before and after'' (\textit{Physics} IV.11, 219b12--13); ``time, then, also is both made continuous by the `now' and divided at it'' (\textit{Physics} IV.12, 220a5). The `now' is at once the principle of continuity --- the moving boundary that links successive durations --- and the principle of division --- the point at which past ends and future begins. ``What is \textit{before} is in time; for we say `before' and `after' with reference to the distance from the `now', and the `now' is the boundary of the past and the future'' (\textit{Physics} IV.14, 223a4--6). Aristotle illuminates the `now's mode of being through the body-carried-along analogy: ``the `now' corresponds to the body that is carried along, as time corresponds to the motion. For it is by means of the body that is carried along that we become aware of the before and after in the motion, and if we regard these as countable we get the `now''' (\textit{Physics} IV.11, 219b22--25). The moving body is always the same body even as it traverses different positions; the `now' is always the same `now' even as it marks different moments. One in substrate, different in being.

Heidegger reads the Aristotelian analysis in \textit{Being and Time} as the most rigorous articulation of what he calls within-time-ness (\textit{Innerzeitigkeit}) --- time as publicly encountered, the time-in-which entities and events occur. Aristotle's analysis is correct as far as it goes: it captures the structure of public, world-time, the time within which the chain's actualizations unfold and which is articulated by the now's bounding and joining function. What the analysis does not thematize is the more primordial structure of temporality (\textit{Zeitlichkeit}) and historicity (\textit{Geschichtlichkeit}) from which \textit{Innerzeitigkeit} is derived. The derivative time of the now-sequence presupposes Dasein's ecstatic temporalizing as the originary source from which the public-time concept is generated through a levelling-off. Heidegger states the relation in his explicit engagement with Aristotle's definition: ````For this is time: that which is counted in the movement which we encounter within the horizon of the earlier and later.'' This definition may seem strange at first glance; but if one defines the existential-ontological horizon from which Aristotle has taken it, one sees that it is `obvious' as it at first seems strange, and has been genuinely derived'' (BT \S 81, H.421; Eng.~473).\footnote{The passage appears in \textit{Being and Time} \S 81 (``Within-time-ness and the genesis of the ordinary conception of time''), in Heidegger's explicit engagement with Aristotle's \textit{Physics} IV.11, 219b1--2 definition of time as the number of motion. Heidegger characterizes the Aristotelian definition as ``genuinely derived'' from the existential-ontological horizon of circumspective concern, while indicating that the analysis leaves unthematized the originary temporality from which \textit{Innerzeitigkeit} derives. Continuing at H.421: ``Ever since Aristotle all discussions of the concept of time have clung in principle to the Aristotelian definitions; that is, in taking time as their theme, they have taken it as it shows itself in circumspective concern'' (Eng.~473). The Aristotelian analysis is, on this reading, the most rigorous philosophical account of the within-time-ness of entities, but the source of that time --- the originary temporality from which \textit{Innerzeitigkeit} is derived --- ``does not become a problem for Aristotle'' (Eng.~473).} The structural-philosophical point is precise. Aristotle has captured the \textit{Innerzeitigkeit}-stratum of time without thematizing its grounding in originary temporality; the section's apparatus, accordingly, deploys what Heidegger calls \textit{Innerzeitigkeit} as the time-stratum within which the chain unfolds, while the more primordial temporality that the chain-actualizations presuppose is held open by the Heideggerian framework. The two registers operate together. Aristotle's analysis is the most precise account available of the time-in-which the actualization chain proceeds; Heidegger's analysis is the most precise account available of the temporality from which that public time derives.

The implication for the section's broader project is direct. The cognition of time as a coherent narrative of past, present, and future --- the temporal manifold within which memory, present perception, and anticipation are unified --- is not delivered by the Aristotelian definition alone. The definition captures countable succession; it does not, on its own, articulate the synthetic activity through which the three temporal modes are knit into a unified phenomenal manifold. That synthetic activity is the work of \textit{phantasia}, which the subsequent sections develop in detail. The Aristotelian motion-time horizon within which $A_0$ stands is the substrate within which \textit{phantasia}'s temporal-synthesizing activity unfolds. Aristotle himself prepares the thesis:
\begin{adjustwidth}{0.5in}{0in}
Further, pleasure is the consciousness through the senses of a certain kind of emotion; but imagination is a feeble sort of sensation, and there will always be in the mind of a man who remembers or expects something the imagination of what he remembers or expects. If this is so, it is clear that memory and expectation also, being accompanied by sensation, may be accompanied by pleasure. Hence it follows that everything pleasant is either present in sensation, remembered, or future and expected, since we perceive present things, remember past ones, and expect future ones. (\textit{Rhetoric} I.11, 1370a27--b1)
\end{adjustwidth}
\textit{Phantasia} is what allows the soul to be moved by absences, possibilities, appearances, and counterfactuals; it is the capacity that articulates the motion-time substrate into the temporal manifold within which past, present, and future co-temporalize for the comporting being.

\subsection{The Soul's Constitutive Role}

Aristotle's analysis of the soul in \textit{De Anima} establishes the soul's constitutive role in the actualization chain. The soul is ``a substance [\textit{ousia}] in the sense of the form [\textit{eidos}, \gk{εἶδος}] of a natural body'' (\textit{De Anima} II.1, 412a19--20); it is the first actuality of a natural body that has life potentially, what makes a body into a living organism capable of perceiving, moving, and thinking. The soul does not exist apart from the body; it is the body's organizing principle, its way of being what it is. Hence the soul is itself an instance of the priority of actuality: it is the form that must be actual before the body's capacities can be exercised.\footnote{Pöggeler reads Heidegger's recovery of Aristotelian soul as the form of a living body as the conceptual scaffold for the entire \textit{Bewegtheit} (``movement'') as being-underway-toward reading of GA 18, which is then partially retracted in Heidegger's later turn: ``When attunement turns toward or away from situational moods, movement and time come into play.... The element of the pathetic or of feeling is relegated to attunement (Befindlichkeit) and mood (Stimmung), in which being-in-the-world suppressed becomes manifest'' (Pöggeler 169). This 1924 reading is, in Pöggeler's interpretation, the originary site of the analyses Heidegger would later consolidate as \textit{Befindlichkeit}/\textit{Stimmung}. Indeed, it is the priority of actuality at the soul-level, as I understand it, that makes the temporal-affective structure of being-in-the-world possible at all.} The point connects directly to the motion-time parallel. Just as motion is the actuality of the potential qua potential, so the soul is the actuality of the body's potential for life; just as time is the numbered aspect of motion, so the rational soul is the numbering principle that articulates temporal flow.\footnote{The structural parallel runs between the soul as first actuality of an organic body (\textit{De Anima} II.1, 412a27--b1: ``the soul is the first actuality [\textit{entelecheia}] of a natural body which has life in potentiality'') and time as the number of motion requiring soul to count (\textit{Physics} IV.11, 219b1--2: ``time is the number of movement in respect of the before and after''; cf.\ \textit{Physics} IV.14, 223a21--26 on time's dependence on soul). Aristotle himself notes the soul-time connection (\textit{Physics} IV.14) but does not extend the actuality-motion analogy explicitly to it; the parallel is supported by the shared \textit{energeia}-structure on both sides --- soul as first \textit{entelecheia}, motion as the \textit{energeia} of the potential \textit{qua} potential, time as the numbered aspect of that motion --- such that the rational soul stands to the temporal manifold as form stands to enmattered actualization.}

Aristotle's three-grade schema of potentiality clarifies the architectural placement of $A_1$: first potentiality is the bare capacity, second potentiality or first actuality is the developed capacity, second actuality is the exercise of that capacity (\textit{De Anima} II.5, 417a21--b2). $A_1$ corresponds to second actuality: the completed perceptual act issuing from $A_0$'s dyadic actuality, in which the developed first-actuality capacities of the sense-organ and the actually-exercising sensible object --- the constituents of $A_0$ --- are jointly engaged. The hierarchical ordering of psychic capacities (\textit{De Anima} II.3) reinforces the iterated structure. Each level presupposes the actualization of those below it: perception presupposes the nutritive capacity, \textit{phantasia} presupposes perception, thought presupposes \textit{phantasia}.

The architectural weight of the motion-time horizon becomes most explicit at the soul's relation to time. ``If nothing but soul, or in soul reason, is qualified to count, it is impossible for there to be time unless there is soul, but only that of which time is an attribute --- i.e., if movement can exist without soul'' (\textit{Physics} IV.14, 223a21--26). The conditional is precise. Motion may exist without soul: natural bodies move whether or not anyone observes them. But time --- the numbered aspect of motion --- requires a counter, and where there is no counter there can be no number counted. The section adopts a stronger reading of the conditional than Aristotle's text strictly compels: time without a counting soul is not merely unnumbered but ontologically incomplete: it is present as countability, structurally real as the before-and-after of motion, yet not fully actualized as time.\footnote{Coope, in \textit{Time for Aristotle: Physics IV.10--14} (Oxford, 2005), Ch.~10 ``Time and the Soul'', defends a mind-dependence reading: ``Once we fully understand Aristotle's view about the way in which time is countable, we should be able to see for ourselves not just that it is mind-dependent but also that it is mind-dependent in a way in which change is not'' (p.~160). Her position: ``the nature of time itself implies that time cannot exist in the absence of ensouled beings'' (p.~162). Broadie defends a contrasting reading on which the before-and-after structure of motion persists in a soulless world, with only time-\textit{as-arithmos} requiring a counter: ``****** UNVERIFIED'' (Broadie, \textit{Nature, Change, and Agency in Aristotle's Physics} (Oxford, 1982); UNVERIFIED --- Broadie 1982 not currently in corpus; user to supply verified verbatim). Aristotle's own doctrine of the co-eternity of motion and time (\textit{Metaphysics} XII.6, 1071b8--11, quoted above) entails that time exists wherever motion exists, since motion's before-and-after structure is intrinsic to it. But his account of latent versus actualized potentiality (\textit{Metaphysics} IX.9, 1051a29--33, quoted above) entails that countable structures present in a configuration are not fully actual until brought into actuality by a thinker: as the geometer's potential constructions exist in the figure but are actualized only when the geometer proves them, so time as the number of motion is structurally present in motion's before-and-after yet is fully actualized as numbered only when the rational soul performs the counting.} It coheres with Aristotle's broader priority-of-actuality commitments: if time is fully actual only as numbered, what makes the numbering possible is constitutive of time's full actuality.

Heidegger radicalizes the soul-as-counter thesis as originary temporality (\textit{Zeitlichkeit}). Originary temporality is not a sequence of nows but the unity of three ecstases (\textit{Ekstasen}): having-been (\textit{Gewesenheit}), present (\textit{Gegenwart}), and projected-futural (\textit{Zukunft}). Each ecstasis is a mode in which Dasein temporalizes itself; the three are co-original, and their unity is the originary phenomenon from which all derivative time-conceptions are drawn. Heidegger states the relation: ``Temporality is the primordial `out-side-of-itself' in and for itself. We therefore call the phenomena of the future, the character of having been, and the Present, the \textquotedblleft ecstases\textquotedblright\ of temporality'' (BT \S 65, H.328--329; Eng.~377). The three ecstases do not arise sequentially: ``temporality does not first arise through a cumulative sequence of the ecstases, but in each case temporalizes itself in their equiprimordiality'' (BT \S 65, H.329; Eng.~378). The everyday now-sequence is derivative: ``the `time' which is accessible to Dasein's common sense is not primordial, but arises rather from authentic temporality'' (BT \S 65, H.329; Eng.~377).\footnote{Heidegger develops the three ecstases as the unitary temporalizing-source: ``Temporality `is' not an entity at all. It is not, but it temporalizes itself'' (BT §65, H.328; Eng.~377); ``its essence is a process of temporalizing in the unity of the ecstases'' (BT §65, H.329; Eng.~377). The future has priority within this equiprimordiality: ``The primary phenomenon of primordial and authentic temporality is the future'' (BT §65, H.329; Eng.~378). The derivation of \textit{Innerzeitigkeit} from originary temporality is developed across BT §§66--68 (H.331--344); the historicity (\textit{Geschichtlichkeit}) thesis builds on the same temporalizing-source: ``As historical, Dasein is possible only by reason of its temporality, and temporality temporalizes itself in the ecstatico-horizonal unity of its raptures'' (BT §74, H.385--386; Eng.~437). Aristotle's soul-as-counter operates at the categorial-structural level on what Heidegger's circumspective concern (BT §§66--68) temporalizes at the existential-ontological level.} The Aristotelian insight that time requires soul as counter is secured at the existential-ontological level by this analysis: Dasein's ecstatic temporality \textit{is} the originary phenomenon from which the public-time of the now-sequence derives. The actualization chain's public time is, accordingly, an articulated stratum derivative from Dasein's temporalizing --- the motion-time horizon and the rational soul that articulates it together constitute the time-stratum within which the chain proceeds.

\subsection{The Threshold of $A_1$}

The threshold of $A_1$ is the architectural pivot at which the section's central thesis enters: \textit{phantasia} is the temporal-synthetic capacity through which a finite, embodied, comporting being cognizes time as a coherent narrative of past, present, and future. Aristotle states the foundation: ``all memory implies time; therefore only animals that perceive time can remember, and the organ by which they perceive time is that by which they remember'' (\textit{De Memoria} 449b28--30); and ``imagination is a feeble sort of sensation'' (\textit{Rhetoric} I.11, 1370a28). \textit{Phantasia} is the capacity that ranges across all three temporal modes --- it presents the present object, retains the absent past object as memory, and projects the not-yet object as anticipation. Without \textit{phantasia}, the soul could be moved only by what is perceptually present; with \textit{phantasia}, the soul can be moved by absences, possibilities, appearances, and counterfactuals, which is to say by the full range of objects required for the actualization chain to operate across memory and anticipation rather than within the immediate `now' alone.\footnote{Frede emphasizes that ``\textit{phantasiai} can thus be separated from their origin, while perceptions cannot, and this means that they can give us a coherent picture of a situation that transcends the immediate perception'' (Frede 285). The temporal range that this project ascribes to \textit{phantasia} --- its capacity to range across past memory, present perception, and future anticipation --- is, on Frede's reading, the cognitive-architectural consequence of the residual-motion structure: once the \textit{phantasma} is detached from its original perceptual source, it can be deployed across the whole temporal horizon the rational soul articulates.} Heidegger's productive imagination radicalizes this Aristotelian capacity at the level of fundamental ontology. In \textit{Kant and the Problem of Metaphysics} (§35) he isolates the productive imagination as the ``root of both stems'' --- sensibility and understanding --- and shows that this root is itself ``rooted in original time'' (Taft trans., §35, p.~141), which is ``the original, threefold-unifying forming of future, past, and present in general'' (Taft trans., §35, p.~137).\footnote{The structural homology between Aristotelian \textit{phantasia} and Heideggerian productive imagination is that both name the synthesizing capacity through which discrete temporal contents --- the present perceptual instant, the past memory, the future anticipation --- are unified into a single temporal manifold. For Aristotle, this synthesizing role is what makes the chain's traversal across temporal modes possible: \textit{phantasia} preserves the perceptual residue as memory, projects the anticipatory image, and presents the current object as a unified phenomenal whole. For Heidegger, productive imagination is ``the root of both stems, sensibility and understanding. As such, it makes possible the original unity of ontological synthesis. This root, however, is rooted in original time. The original ground which becomes manifest in the ground-laying is time'' (Taft trans., §35, p.~141). The two analyses operate at different levels --- Aristotle at the level of psychic capacity, Heidegger at the level of transcendental synthesis grounding fundamental ontology --- but converge on the structural claim that temporal unity is a productive achievement of a single synthesizing capacity, not a given of consciousness.} The two analyses converge on the same structural claim: cognition of time as a unified manifold of past, present, and future is the work of a single synthesizing capacity, without which the comporting being could be moved only by what is immediately present. The full development of the \textit{phantasia}-thesis belongs to the subsequent sections; what the motion-time horizon together with $A_0$ establishes is the ontological substrate within which \textit{phantasia}'s temporal-synthesizing activity becomes intelligible, and it is by virtue of this synthesizing activity that the chain spans past, present, and future as a coherent narrative manifold.

With the phantasia-thesis named, the threshold of $A_1$ can be stated with precision. The motion-time horizon, together with $A_0$'s dyadic actuality, licenses the preconditions of $A_1$: natural beings consisting of sensible objects exercising their sensible powers, sense-organs holding their first actuality as capacities, and rational souls capable of articulating temporal succession. Motion being actual is what allows sensible objects to actually exercise their sensible powers --- the color visible, the sound audible, the warm warming. Motion being actual is what allows the medium to transmit the sensible form from object to organ. Time being actual is what allows the temporal continuum within which the transmission unfolds; and time becoming fully actual with the arrival of rational soul is what allows the sequential articulation of perceptual, imaginative, and cognitive acts the subsequent sections will trace. $A_1$, completed perception, is the second actuality in which all four preconditions jointly engage: rational souls perceiving sensible objects actually exercising their sensible powers, sense-organs holding first actuality as capacity, and a temporal-narrative manifold within which the actualization chain proceeds. Each precondition depends on $A_0$, and $A_1$ itself will become the unmoved originator of the phantastic motion as the chain unfolds. The chain begins.

\end{document}